% Source: https://artofproblemsolving.com/community/c3h2242889_view_question_mathcounts_state_team_9_2006
% Stable identity: aops-topic-2242889-q01
A rectangle using all $n$ tiles is an unordered pair $\{a,b\}$ of positive integers with $ab=n$. If $n$ is not a perfect square, the $d(n)$ divisors of $n$ pair up, giving $\dfrac{d(n)}{2}$ rectangles; if $n$ is a square, there are $\dfrac{d(n)+1}{2}$ (the square counts once). Ten rectangles therefore need $d(n)=20$ with $n$ not a square, or $d(n)=19$ with $n$ a square.

If $d(n)=19$, then $n=p^{18}$ for a prime $p$, so $n\ge2^{18}=262144$.

If $d(n)=20$ and $n=p_1^{e_1}p_2^{e_2}\cdots$, then $(e_1+1)(e_2+1)\cdots=20$, with the factorisations $20$, $10\cdot2$, $5\cdot4$, $5\cdot2\cdot2$. Putting the largest exponents on the smallest primes, the smallest numbers of these forms are
\[2^{19},\qquad 2^9\cdot3=1536,\qquad 2^4\cdot3^3=432,\qquad 2^4\cdot3\cdot5=240 .\]
The least is $\boxed{240}$, which is not a square and indeed has $5\cdot2\cdot2=20$ divisors: its rectangles are $1\times240$, $2\times120$, $3\times80$, $4\times60$, $5\times48$, $6\times40$, $8\times30$, $10\times24$, $12\times20$, $15\times16$.
